\documentclass[border=10pt]{standalone}
\usepackage[UTF8]{ctex}
\usepackage{tikz}
\usepackage{amsmath,amssymb}
\usetikzlibrary{arrows.meta,calc}
\definecolor{cBlue}{RGB}{37,99,235}
\definecolor{cOrange}{RGB}{234,88,12}
\definecolor{cGray}{RGB}{100,116,139}
\definecolor{cGreen}{RGB}{22,163,74}
\definecolor{cRed}{RGB}{220,38,38}
\begin{document}
\begin{tikzpicture}[>=Stealth,line width=0.9pt]
  % M with two overlapping patches
  \draw[cBlue,line width=1.2pt,fill=cBlue!6]
    (-2.2,2.2) to[out=20,in=160] (2.2,2.2) to[out=-20,in=20] (1.8,0.7)
    to[out=200,in=-20] (-1.8,0.7) to[out=160,in=-160] (-2.2,2.2);
  \node[cBlue] at (2.0,2.35) {$M$};
  % patch U (orange) and V (green), overlapping
  \draw[cOrange,line width=1pt] (-0.65,1.55) ellipse (0.95 and 0.55);
  \draw[cGreen,line width=1pt] (0.65,1.55) ellipse (0.95 and 0.55);
  \node[cOrange] at (-1.35,1.55) {\scriptsize $U$};
  \node[cGreen] at (1.35,1.55) {\scriptsize $V$};
  \begin{scope}
    \clip (-0.65,1.55) ellipse (0.95 and 0.55);
    \fill[cRed!18] (0.65,1.55) ellipse (0.95 and 0.55);
  \end{scope}
  \node[cRed] at (0,1.95) {\scriptsize $U\cap V$};
  % two coordinate charts below
  \draw[cOrange!70,fill=cOrange!8] (-3.4,-1.2) rectangle (-1.6,0.1);
  \draw[cGray!45,step=0.3] (-3.4,-1.2) grid (-1.6,0.1);
  \node[cOrange] at (-2.5,-1.45) {\scriptsize $\varphi_U(U)\subset\mathbb{R}^n$};
  \fill[cRed!18] (-2.0,-0.55) ellipse (0.32 and 0.42);
  \draw[cGreen!70,fill=cGreen!8] (1.6,-1.2) rectangle (3.4,0.1);
  \draw[cGray!45,step=0.3] (1.6,-1.2) grid (3.4,0.1);
  \node[cGreen] at (2.5,-1.45) {\scriptsize $\varphi_V(V)\subset\mathbb{R}^n$};
  \fill[cRed!18] (2.0,-0.55) ellipse (0.32 and 0.42);
  % chart maps
  \draw[->,cOrange,line width=1pt] (-1.1,1.05) to[bend right=12] (-2.5,0.2);
  \node[cOrange] at (-2.05,0.95) {$\varphi_U$};
  \draw[->,cGreen,line width=1pt] (1.1,1.05) to[bend left=12] (2.5,0.2);
  \node[cGreen] at (2.05,0.95) {$\varphi_V$};
  % transition map between the two coordinate images
  \draw[->,cRed,line width=1.2pt] (-1.6,-0.55) to[bend right=22] (1.6,-0.55);
  \node[cRed,align=center] at (0,-0.95)
    {$\varphi_V\!\circ\varphi_U^{-1}$\\[-1pt]\footnotesize（换坐标 / 转移映射）};
  \node[align=center,text=cGray] at (0,-2.1)
    {\footnotesize \textbf{光滑流形}：两套坐标卡在重叠处的\textbf{换坐标映射}（$\mathbb{R}^n\!\to\!\mathbb{R}^n$）光滑。\\
     \footnotesize “光滑”说的是这些转移映射，而不是 $M$ 本身——$M$ 只是个拓扑空间。};
\end{tikzpicture}
\end{document}
